\begin{frame}{Il modello}
\phantomsection\label{il-modello}
\(Y=C+I+\hat{G}\)

\(C=\bar{C}+c(Y-\hat{T})\)

\(I=\bar{I}-bR\)

\(\frac{\hat{M}}{\bar{P}}=kY-hR\)

\vskip10mm

Dove \(Y\), \(C\), \(I\) e \(R\) sono le variabili endogene;

\(\bar{C}, \bar{I}, \bar{P}\) variabili esogene non controllabili;

\(\hat{G}, \hat{T}, \hat{M}\), variabili esogene controllabili dalle
autorità;

\(c, b, k, h\) sono parametri.
\end{frame}

\begin{frame}{forma matriciale}
\phantomsection\label{forma-matriciale}
Riorganizziamo in modo che le colonne riportino \(Y\), \(C\), \(I\) e
\(R\)

\(1\times Y-1\times C-1\times I+0\times R=\hat{G}\)

\hskip-2.5mm\(-c\times Y+1\times C+0\times I+0\times R=\bar{C}-c\hat{T}\)

\(0\times Y+0\times C+1\times I+b\times R=\bar{I}\)

\(k\times Y+0\times C+0\times I-h\times R=\frac{\hat{M}}{\bar{P}}\)

In forma matriciale

\[\mathbf{A}x=\mathbf{b}\]

dove \[
\mathbf{A}=\left[
\begin{array}{cccc}
1&-1&-1&0\\
-c&1&0&0\\
0&0&1&b\\
k&0&0&-h
\end{array}
\right]
\qquad 
\mathbf{x}=\left[
\begin{array}{c}
Y\\
C\\
I\\
R
\end{array}
\right]
\qquad 
\mathbf{b}=\left[
\begin{array}{c}
\hat{G}\\
\bar{C}-c\hat{T}\\
\bar{I}\\
\frac{\hat{M}}{\bar{P}}
\end{array}
\right]
\]
\end{frame}

\begin{frame}{La fattorizzazione \(\mathbf{LU}\)}
\phantomsection\label{la-fattorizzazione-mathbflu}
\begin{itemize}
\item
  Una delle tecniche utilizzate per risolvere sistemi lineari con
  matrice \(\mathbf{A}\) quadrata è la fattorizzazione \(\mathbf{LU}\).
\item
  In questa fattorizzazione, vengono identificate due matrici:

  \begin{itemize}
   
  \item
    \(\mathbf{L}\) che è triangolare inferiore (Lower triangular)
  \item
    \(\mathbf{U}\) che è triangolare superiore (Upper triangular)
  \end{itemize}
\end{itemize}

legate ad \(\mathbf{A}\) dalla seguente relazione:

\[\mathbf{LU}=A\]
\end{frame}

\begin{frame}{Soluzione del sistema}
\phantomsection\label{soluzione-del-sistema}
Il sistema \[\mathbf{A}x=\mathbf{b}\] può dunque essere scritto come
\[\mathbf{LU}x=\mathbf{b}\]

La soluzione è ottenuta in due passaggi:

prima risolviamo

\(\mathbf{Ly=b}\)

rispetto a \(\textbf{y}\), dove \(\mathbf{y=Ux}\)

e poi risolviamo il sistema

\(\mathbf{Ux=y}\)

ottenedo così la soluzione per \(\mathbf{x}\).
\end{frame}

\begin{frame}{Esistenza della fattorizzazione \(\mathbf{LU}\)}
\phantomsection\label{esistenza-della-fattorizzazione-mathbflu}
La fattorizzazione \(\mathbf{LU}\) potrebbe non esistere anche se
\(\mathbf{A}\) è non singolare.

\vskip5mm

Proposizione: Data una matrice \(\mathbf{A}\in R^{n\times n}\), la sua
fattorizzazione \(\mathbf{LU}\) esiste ed è unica se e solo se le
sottomatrici principali \(\mathbf{A_i}\) di \(\mathbf{A}\) di ordine
\(i = 1,...,n-1\) (cioè quelle ottenute limitando \(\mathbf{A}\) alle
sole prime \(i\) righe e colonne) sono non singolari.

\vskip5mm

Zhang, F.: Matrix Theory. Universitext. Springer, New York (1999)
paragrafo 3.2
\end{frame}

\begin{frame}{Casi noti}
\phantomsection\label{casi-noti}
La fattorizzazione \(\mathbf{LU}\) esiste se la matrice è

\begin{itemize}
\item
  a dominanza diagolale stretta
\item
  reale simmetrica e definita positiva
\item
  complessa e definita positiva
\end{itemize}

Nel secondo caso esiste la fattorizzazione speciale nota come
fattorizzazione di Cholesky:

\[A=\mathbf{R^TR}\] dove \(\mathbf{R}\) è una matrice triangolare
superiore con elementi positivi sulla diagonale.

In questo caso dobbiamo identificare solo una matrice e il numero di
operazioni necessarie si dimezza rispetto alla fattorizzazione
\(\mathbf{LU}\).
\end{frame}

\begin{frame}{Richiami sulle matrici (dominanza diagonale)}
\phantomsection\label{richiami-sulle-matrici-dominanza-diagonale}
\begin{itemize}
 
\item
  Una matrice è a dominanza diagonale per righe se per ogni riga, il
  valore assoluto dell'elemento che appartiene alla diagonale è \(\ge\)
  al valore della somma del valore assoluto degli elementi fuori
  diagonale. \[|a_{ii}|\ge \sum_{j\ne i} |a_{ij}|\] La dominanza è
  stretta se la disuguaglianza è stretta (\(>\))
\item
  La definizione di dominanza diagonale per colonna segue la stessa
  logica di quella per riga.
\end{itemize}
\end{frame}

\begin{frame}{Richiami sulle matrici (definita positiva)}
\phantomsection\label{richiami-sulle-matrici-definita-positiva}
\begin{itemize}
 
\item
  Una matrice è definita positiva se \[x^TAx>0\] \(\forall x\ne0\). Se
  la disuguaglinza non è stretta la matrice si dice semi definita
  positiva.
\end{itemize}

\vskip1cm

Si noti che \(x^T\) è una matrice \(1\times n\) e \(A\) una matrice
\(n\times n\). Il prodotto \(x^TA\) sarà dunque una matrice
\(1\times n\).

\(x\) è una matrice \(n\times 1\).

Moltiplicando dunque una matrice \(1\times n\) (\(x^TA\)) per una
matrice \(n\times 1\) (\(x\)) si ottiene una matrice \(1\times 1\) che è
un singolo numero di cui può essere valutato il segno.
\end{frame}

\begin{frame}{Come ottenere le matrice L e U (cenno e rinvio)}
\phantomsection\label{come-ottenere-le-matrice-l-e-u-cenno-e-rinvio}
Imponiamo l'uguaglianza \[
\left[
\begin{array}{ccc}
l_{11}&0&0\\
l_{21}&l_{22}&0\\
l_{31}&l_{32}&l_{33}
\end{array}
\right]
\left[
\begin{array}{ccc}
u_{11}&u_{12}&u_{13}\\
0&u_{22}&u_{23}\\
0&0&u_{33}
\end{array}
\right]=
\left[
\begin{array}{ccc}
a_{11}&a_{12}&a_{13}\\
a_{21}&a_{22}&a_{23}\\
a_{31}&a_{32}&a_{33}
\end{array}
\right]
\] dove \(l_{ij}\) e \(u_{ij}\), sono le incognite. Si noti che abbiamo
\(x\times n +n\) incognite.

effettuando i prodoti riga per colonna si ottengono \(n\times n\)
equazioni.

Avendo più incognite che equazioni semplifichiamo mettendo degli uno
nella diagonale di \(\mathbf{L}\)
\end{frame}

\begin{frame}{}
\phantomsection\label{section}
\[
\left[
\begin{array}{ccc}
1&0&0\\
l_{21}&1&0\\
l_{31}&l_{32}&1
\end{array}
\right]
\left[
\begin{array}{ccc}
u_{11}&u_{12}&u_{13}\\
0&u_{22}&u_{23}\\
0&0&u_{33}
\end{array}
\right]=
\left[
\begin{array}{ccc}
a_{11}&a_{12}&a_{13}\\
a_{21}&a_{22}&a_{23}\\
a_{31}&a_{32}&a_{33}
\end{array}
\right]
\] impostiamo ora il sistema di equazioni:

\vskip5mm

\(u_{11}=a_{11}\) (prima riga per prima colonna)

\(\cdots\)

\(\cdots\)

\(l_{31}u_{13}+l_{32}u_{23}+u_{33}=a33\) (ultima riga per ultima
colonna)

\vskip5mm

le prime equazioni sono molto semplici da risolvere, e pertanto il
carico computazionale si riduce.
\end{frame}

\begin{frame}{Rinvio}
\phantomsection\label{rinvio}
\begin{itemize}
\item
  Il lettore interessato può trovare l'algoritmo per ottenere la
  scomposizione LU a pag. 161 del libro Calcolo Scientifico di Saleri e
  altri.
\item
  A pagine 166 dello stesso testo si trova l'algoritmo per il calcolo
  della fattorizzazione di Cholesky.
\item
  \(\textnormal{\`E}\) possibile giungere alla scomposizione LU anche
  nei casi in cui non siano verificate le condizioni esposte alla
  dipositiva 6. Il meccanismo prevede lo scambio di righe della matrice
  scegliendo quelle che hanno elementi con valore più alto in certe
  posizioni (elementi detti pivot). Si rinvia a pagina 169 del libro
  citato per l'esposizione della tecnica del pivoting.
\end{itemize}
\end{frame}

\begin{frame}{}
\phantomsection\label{section-1}
Abbiamo in precedenza evidenziato i problemi di pesantezza
computazionale del calcolo del determinante (e di conseguenza del
calcolo dell'inversa e della soluzione dei sistemi lineari tramite il
metodo di Kramer).

\vskip5mm

Poiché se \(\mathbf{C=AB}\) allora
\(det(\mathbf{C})=det(\mathbf{A})det(\mathbf{B})\)

\vskip5mm

nel nostro caso avremo
\(det(\mathbf{A})=det(\mathbf{L})det(\mathbf{U})\)

\vskip5mm

ma, ricordiamo che, il determinante di matrici triangolari è il prodotto
dei termini in diagonale. Considerando che \(\mathbf{L}\) ha tutti 1 in
diagonale abbimo \[det(\mathbf{A})=\prod_1^n u_{ii}\] che si calcola
molto agevolmente.
\end{frame}

\begin{frame}{comandi}
\phantomsection\label{comandi}
R (pacchetto Matrix)

lu(A)

\vskip5mm

Octave

lu(a)

\vskip5mm

python

scipy.linalg.lu(a, permute\_l=False, overwrite\_a=False,
check\_finite=True)

\vskip5mm

julia

LinearAlgebra.lu(A, pivot = RowMaximum(); check = true)
\end{frame}

\begin{frame}{Algoritmo}
\phantomsection\label{algoritmo}
\begin{itemize}
 
\item
  assegnare matrice \(\mathbf{A}\) e vettore \(\mathbf{d}\)
\item
  ottenere le matrici \(\mathbf{L}\) e \(\mathbf{U}\) utilizzando la
  funzione messa a disposizione dal software
\item
  ottenere \(\mathbf{y}\) utilizzando \(\mathbf{L}\) e \(\mathbf{d}\)
\item
  ottenere \(\mathbf{x}\) utilizzando \(\mathbf{U}\) e \(\mathbf{y}\)
\end{itemize}
\end{frame}
